\input{Iteration-V2/Chapter4/chapter-04-12}

\subsection{Revision 30 continuity consolidation}

The submission-length edition retains the established evidence-bounded architecture while consolidating post-Revision-14 extensions that otherwise repeated change history across successive drafts. These add source-byte verification, signed anti-rollback trust-store envelopes, witness quorum before reviewer reassurance, explicit empty production authority, and a non-collecting decision pause. Their common rule is unchanged: engineering controls cannot create legal authority, reviewer qualification, participant evidence, or a compliance verdict. The original revision sources remain preserved in the repository.

\subsection{Revision 30 implementation: declarative mapping and decoupled insertion}

The implementation separates three responsibilities. The regulation pack stores the declarative temporal profiles alongside its ordinary governance metadata. A dedicated conversion module validates profile identifiers, supported observation types, positive safe counts, legal references, rationale, and a bounded non-zero window before compiling profiles into the generic evaluator's representation. The evaluator itself accepts only compiled rules and package-scoped observations; it imports no GDPR-specific list and therefore remains reusable when another regulation pack is mounted. The governance validator invokes the mapping validator, so a malformed profile fails closed before it can alter a new review.

This separation preserves linkage without hard coupling. The pack is the owner of legal meaning; the mapping compiler is the explicit translation boundary; the temporal evaluator owns multiset matching, canonical ordering, receipt hashing, and bounded-ledger treatment; and the compliance engine owns the existing finding lifecycle. Swapping the mounted pack changes the valid windows and combinations through the same interface. A small research-baseline pack intentionally contains a different sensor-fusion profile, demonstrating that the mapping is data-driven rather than a hidden GDPR fallback.

The Android bridge enriches ordinary audit evidence with normalised observation events when an available bridge API/action can establish a type, channel, source, and non-zero count. It does not infer absent calls, and a zero frequency produces no temporal observation. The TypeScript compliance engine keeps a bounded per-package ledger, validates the active mapping, deduplicates matched evidence, and emits a neutral temporal-cooccurrence finding containing the rule identifier, window, requirements, observed receipt, legal references, rationale, and source classification. The interface exposes these fields as an evidence card so that the reviewer can see the window and combination rather than receiving an opaque severity label.

The feature remains insertable rather than entangled with permission interception. It invokes the same notification and evidence path used by the existing review engine; it neither denies an API request nor terminates an application. This is a conscious constraint: an order-independent cluster of observations can justify a question for review, but without purpose, necessity, data-flow, consent, retention, and lawful-basis evidence it cannot safely justify enforcement.
